On considère un quadrillage. Pour aller du point $A$ au point $B$ par un
chemin de longueur minimale, on est uniquement autorisé à effectuer :
\begin{itemize}
  \item un déplacement vers la droite, noté $D$ ;
  \item un déplacement vers le haut, noté $H$.
\end{itemize}
Voici un exemple de chemin comportant $4$ déplacements vers la droite et $3$
déplacements vers le haut.

À chaque chemin, on associe le mot formé par la suite de ses déplacements.

Dans l'exemple représenté,
\[
\text{chemin} \longleftrightarrow DDHDHHD.
\]
\begin{enumerate}
  \item Combien de lettres $D$ et de lettres $H$ possède le mot associé au
  chemin représenté ?
  \item Montrer qu'un chemin minimal de $A$ vers $B$ correspond à un unique
  mot contenant $4$ lettres $D$ et $3$ lettres $H$.
  \item Réciproquement, montrer que tout mot contenant $4$ lettres $D$ et $3$
  lettres $H$ définit un unique chemin minimal.
  \item En déduire que le nombre de chemins représentés sur cette grille est
  \[
  \binom{7}{4} = \binom{7}{3}.
  \]
  \item On considère maintenant un quadrillage nécessitant $p$ déplacements
  vers la droite et $q$ déplacements vers le haut. Montrer que les chemins
  minimaux sont en bijection avec les mots de longueur $p + q$ contenant
  exactement $p$ lettres $D$.
  \item En déduire que le nombre de chemins minimaux est
  \[
  \binom{p+q}{p} = \binom{p+q}{q}.
  \]
\end{enumerate}
